\sloppy
\section{Κεφάλαιο 2: Θεωρητικό Υπόβαθρο (Theoretical Background)}

Στο παρόν κεφάλαιο παρουσιάζεται το θεωρητικό πλαίσιο πάνω στο οποίο εδράζονται όλα τα πειράματα της παρούσας διπλωματικής.

\subsection{2.1 Μηχανική Μάθηση σε Βιοϊατρικά Δεδομένα Μικρού Μεγέθους}

Η Μηχανική Μάθηση (Machine Learning - ML), υποπεδίο της Τεχνητής Νοημοσύνης, εστιάζει στην ανάπτυξη αλγορίθμων που μαθαίνουν μοτίβα από δεδομένα και κάνουν προβλέψεις σε άγνωστα δείγματα. Στον χώρο της Υπολογιστικής Βιολογίας και της Ιατρικής Ακριβείας (Precision Medicine), η Επιβλεπόμενη Μάθηση (Supervised Learning) κυριαρχεί. Σε αυτό το παράδειγμα, ο αλγόριθμος εκπαιδεύεται σε ένα σύνολο ζευγών εισόδου-εξόδου $\{(x_1, y_1), \ldots, (x_N, y_N)\}$, με στόχο να προσεγγίσει τη συνάρτηση απεικόνισης $f: X \to Y$ που ελαχιστοποιεί το σφάλμα πρόβλεψης. 

Ωστόσο, η εφαρμογή ML σε βιοϊατρικά δεδομένα, όπως τα πρωτεομικά της παρούσας μελέτης, συνοδεύεται από δομικές προκλήσεις που είναι εγγενώς διαφορετικές από αυτές μιας τυπικής εφαρμογής σε μεγάλα datasets υπολογιστικής όρασης ή επεξεργασίας φυσικής γλώσσας.

Η πρώτη και κρισιμότερη πρόκληση είναι η \textbf{Κατάρα της Διαστασιμότητας} (Curse of Dimensionality), όπως τη διατύπωσε ο Richard Bellman το 1957. Καθώς αυξάνεται ο αριθμός των διαστάσεων (χαρακτηριστικών/πρωτεϊνών), ο όγκος του χώρου αυξάνεται εκθετικά, καθιστώντας τα διαθέσιμα δεδομένα αραιά (sparse). Για να επιτευχθεί στατιστική σημαντικότητα, ο αριθμός των δειγμάτων $N$ θα έπρεπε να αυξάνεται εκθετικά με τις διαστάσεις $P$. Στην περίπτωσή μας, όπου $P \approx 500$ (πρωτεΐνες) και $N < 100$ (ασθενείς), υπάρχει μεγάλος κίνδυνος το μοντέλο να βρει τυχαίες συσχετίσεις (spurious correlations), οδηγώντας σε υπερ-προσαρμογή (overfitting) [25], [30]. 

Η αντιμετώπιση απαιτεί συνδυασμό τεχνικών μείωσης διαστάσεων, regularization και επιλογής χαρακτηριστικών.

Η δεύτερη πρόκληση είναι η \textbf{Ανισορροπία Κλάσεων} (Class Imbalance). Στα ιατρικά σύνολα δεδομένων, οι ασθενείς με τη νόσο (θετικά δείγματα) είναι συχνά μειοψηφία. Ένας ταξινομητής που εκπαιδεύεται ελαχιστοποιώντας απλώς το συνολικό σφάλμα (accuracy), τείνει να μεροληπτεί υπέρ της πλειοψηφούσας κλάσης (Healthy). Αυτό αντιμετωπίζεται μέσω τεχνικών όπως η Σταθμισμένη Συνάρτηση Κόστους (Weighted Loss Function) και η Στρωματοποιημένη Δειγματοληψία (Stratified Sampling) [44], [45].

Σε αυτό το γενικό πλαίσιο, δύο μεγάλες οικογένειες αλγορίθμων κυριαρχούν στη βιβλιογραφία: τα Νευρωνικά Δίκτυα (Neural Networks), τα οποία έχουν αποδειχθεί ιδιαίτερα ισχυρά σε προβλήματα με μεγάλα δεδομένα και πολύπλοκη δομή, και οι Μέθοδοι Συνόλου (Ensemble Methods) βασισμένες σε δέντρα αποφάσεων, οι οποίες είναι θεωρητικά καταλληλότερες για μικρότερα, δομημένα (tabular) datasets όπως αυτό της παρούσας μελέτης. Στις ενότητες που ακολουθούν αναπτύσσονται και οι δύο, στη σειρά που εφαρμόζονται και στα πειράματα.

\subsection{2.2 Νευρωνικά Δίκτυα για Ακολουθιακά Δεδομένα}

Τα Τεχνητά Νευρωνικά Δίκτυα (Artificial Neural Networks - ANNs) είναι μαθηματικά μοντέλα εμπνευσμένα από τη βιολογική αρχιτεκτονική του εγκεφάλου. Η θεμελιώδης μονάδα είναι ο τεχνητός νευρώνας (perceptron). Ένας νευρώνας $k$ δέχεται ένα διάνυσμα εισόδου $\mathbf{x} = [x_1, x_2, \ldots, x_m]^T$, το οποίο πολλαπλασιάζεται με ένα διάνυσμα βαρών $\mathbf{w}_k = [w_{k1}, w_{k2}, \ldots, w_{km}]^T$. Σε αυτό προστίθεται μια σταθερά πόλωσης (bias) $b_k$, και το αποτέλεσμα περνά από μια μη γραμμική συνάρτηση ενεργοποίησης $\phi(\cdot)$:
$$y_k = \phi \left( \sum_{j=1}^{m} w_{kj} x_j + b_k \right) = \phi(\mathbf{w}_k^T \mathbf{x} + b_k)$$
Οι μη γραμμικές συναρτήσεις ενεργοποίησης είναι κρίσιμες, διότι χωρίς αυτές, ένα δίκτυο πολλαπλών επιπέδων θα κατέρρεε σε ένα απλό γραμμικό μοντέλο. Στις πιο διαδεδομένες ανήκουν η \textbf{ReLU} $f(x) = \max(0, x)$ (η πιο δημοφιλής λόγω ταχείας σύγκλισης και αποφυγής vanishing gradient στα θετικά πεδία), η \textbf{Sigmoid} $\sigma(x) = 1 / (1 + e^{-x})$ (χρησιμοποιείται στις πύλες των LSTM επειδή παράγει τιμές στο $(0,1)$, ερμηνεύσιμες ως πιθανότητες), και η \textbf{tanh} $\tanh(x) = (e^x - e^{-x})/(e^x + e^{-x})$ (παράγει τιμές στο $(-1,1)$, κεντράροντας τα δεδομένα στο μηδέν, και χρησιμοποιείται ευρέως στα RNNs).


Σε πολλά ιατρικά προβλήματα, η πληροφορία δεν βρίσκεται μόνο στα μεμονωμένα χαρακτηριστικά αλλά και στη χρονική τους αλληλουχία. Τα \textbf{Αναδρομικά Νευρωνικά Δίκτυα (Recurrent Neural Networks - RNNs)}, που εισήχθησαν στη δεκαετία του 1980 (Rumelhart et al., 1986), σχεδιάστηκαν για να χειρίζονται ακολουθιακά δεδομένα. Η βασική τους καινοτομία είναι η ύπαρξη βρόχων ανάδρασης, που επιτρέπουν στην πληροφορία να παραμένει στο δίκτυο. Σε κάθε χρονική στιγμή $t$, το RNN δέχεται την τρέχουσα είσοδο $x_t$ και την κρυφή κατάσταση (hidden state) $h_{t-1}$ από το προηγούμενο βήμα, η οποία λειτουργεί ως η ``μνήμη'' του δικτύου. Η ενημέρωση της κρυφής κατάστασης περιγράφεται από την εξίσωση:
$$h_t = \tanh(W_{hh} h_{t-1} + W_{xh} x_t + \mathbf{b}_h)$$
όπου $W_{hh} \in \mathbb{R}^{d_h \times d_h}$ είναι ο πίνακας βαρών που ελέγχει πώς η προηγούμενη μνήμη επηρεάζει την τρέχουσα (Recurrent weights), $W_{xh} \in \mathbb{R}^{d_h \times d_x}$ ο πίνακας που ελέγχει την επίδραση της νέας εισόδου (Input weights), $\mathbf{b}_h$ το διάνυσμα πόλωσης, και $\tanh$ η συνάρτηση ενεργοποίησης. Τα RNNs έχουν χρησιμοποιηθεί ευρέως στην Επεξεργασία Φυσικής Γλώσσας (NLP) και τη γονιδιωματική. Πάσχουν όμως από το πρόβλημα της \textbf{Εξαφάνισης της Κλίσης} (Vanishing Gradient): κατά την εκπαίδευση με Backpropagation Through Time (BPTT), οι κλίσεις (gradients) πολλαπλασιάζονται διαδοχικά με τον πίνακα $W_{hh}$. Αν οι ιδιοτιμές του $W_{hh}$ είναι μικρότερες του 1, το gradient τείνει εκθετικά στο μηδέν, καθιστώντας αδύνατη τη μάθηση μακροχρόνιων εξαρτήσεων (long-term dependencies).

Για να αντιμετωπιστεί αυτό το πρόβλημα μνήμης των RNNs, οι Hochreiter και Schmidhuber (1997) παρουσίασαν τα \textbf{Long Short-Term Memory} δίκτυα (LSTMs). Η κεντρική ιδέα είναι η εισαγωγή μιας ``λεωφόρου πληροφορίας'', της Κατάστασης Κελιού (Cell State, $C_t$), η οποία επιτρέπει στην πληροφορία να ρέει σχεδόν ανεμπόδιστα στον χρόνο. Η ροή ελέγχεται από δομές που ονομάζονται \textbf{πύλες} (gates). Συγκεκριμένα, το LSTM χρησιμοποιεί τρεις πύλες σε κάθε χρονικό βήμα. Η \emph{Forget Gate} αποφασίζει ποια πληροφορία από το παρελθόν ($C_{t-1}$) είναι άχρηστη και πρέπει να διαγραφεί, μέσω της $f_t = \sigma(W_f \cdot [h_{t-1}, x_t] + b_f)$. Αν $f_t \approx 0$, η μνήμη διαγράφεται· αν $f_t \approx 1$, διατηρείται. Η \emph{Input Gate} αποφασίζει ποια νέα πληροφορία είναι σημαντική για να αποθηκευτεί: $i_t = \sigma(W_i \cdot [h_{t-1}, x_t] + b_i)$, ενώ μια υποψήφια μνήμη δημιουργείται ως $\tilde{C}_t = \tanh(W_C \cdot [h_{t-1}, x_t] + b_C)$. Η ενημέρωση της κατάστασης κελιού συνδυάζει την παλιά μνήμη (αφού φιλτραριστεί από το forget gate) με τη νέα (αφού φιλτραριστεί από το input gate):
$$C_t = f_t \odot C_{t-1} + i_t \odot \tilde{C}_t$$
όπου $\odot$ είναι το γινόμενο κατά στοιχείο (Hadamard product). Η πρόσθεση είναι κρίσιμη: επιτρέπει στο gradient να ρέει χωρίς να πολλαπλασιάζεται συνεχώς, λύνοντας το πρόβλημα του Vanishing Gradient. Τέλος, η \emph{Output Gate} καθορίζει τι θα δει το επόμενο επίπεδο: $o_t = \sigma(W_o \cdot [h_{t-1}, x_t] + b_o)$, και η κρυφή κατάσταση εξάγεται ως $h_t = o_t \odot \tanh(C_t)$. Στην παρούσα μελέτη, η χρήση LSTM συγκρίνεται με το Vanilla RNN ακριβώς για να διαπιστωθεί αν η αυξημένη πολυπλοκότητα (περίπου τετραπλάσιες παράμετροι ανά κελί) δικαιολογείται σε μικρές χρονοσειρές ($T=2$).

Πέρα από την επεξεργασία της χρονικής διαδοχής, ένα νευρωνικό μοντέλο μπορεί να ενισχυθεί με \textbf{Μηχανισμούς Προσοχής} (Attention Mechanisms), εμπνευσμένους από την ανθρώπινη οπτική αντίληψη. Ο μηχανισμός αυτός επιτρέπει στο δίκτυο να εστιάζει επιλεκτικά σε συγκεκριμένα τμήματα της εισόδου. Αρχικά αναπτύχθηκε για τη Μηχανική Μετάφραση (Bahdanau et al., 2014) για να λύσει το πρόβλημα της συμπίεσης μεγάλων προτάσεων σε ένα σταθερό διάνυσμα. Σήμερα, με την έλευση των Transformers (Vaswani et al., 2017), αποτελεί τη βάση των σύγχρονων μοντέλων AI (π.χ. GPT, BERT). Στην ιατρική, ένα μοντέλο ``Black Box'' που δίνει απλώς μια πρόβλεψη (``Καρκίνος: 90\%'') δεν είναι αρκετό, οι κλινικοί απαιτούν να γνωρίζουν το ``Γιατί;''. Ο Attention γεφυρώνει αυτό το χάσμα, αποδίδοντας ένα ``βάρος'' σε κάθε χρονικό βήμα ή χαρακτηριστικό και υποδεικνύοντας πού ``κοίταξε'' το μοντέλο [17], [19]. Αν το βάρος στο $T_1$ είναι υψηλό, η εξέλιξη είναι κρίσιμη· αν το βάρος σε μια πρωτεΐνη είναι υψηλό, αυτή αποτελεί πιθανό βιοδείκτη. Στην εργασία αυτή χρησιμοποιείται ο \emph{προσθετικός μηχανισμός προσοχής} (Additive Attention) του Bahdanau. Έστω $h_1, \ldots, h_T$ οι κρυφές καταστάσεις του RNN. Υπολογίζεται πρώτα ένα alignment score μέσω ενός μικρού νευρωνικού δικτύου, $e_t = \mathbf{v}_a^T \tanh(W_a h_t + b_a)$, όπου $W_a, b_a, \mathbf{v}_a$ είναι εκπαιδεύσιμες παράμετροι. Στη συνέχεια εφαρμόζεται Softmax για να μετατραπούν τα scores σε κατανομή πιθανοτήτων, $\alpha_t = \exp(e_t) / \sum_k \exp(e_k)$, και τέλος υπολογίζεται το context vector ως σταθμισμένο άθροισμα, $c = \sum_t \alpha_t h_t$, το οποίο τροφοδοτείται στον τελικό ταξινομητή.

Τέλος, στο ίδιο νευρωνικό πλαίσιο εντάσσονται και τα \textbf{Συνελικτικά Δίκτυα μιας Διάστασης} (1D-CNNs). Ενώ τα κλασικά Convolutional Neural Networks είναι διάσημα για την επεξεργασία εικόνας (2D), η μονοδιάστατη εκδοχή τους είναι εξαιρετικά ισχυρή στην ανάλυση ακολουθιακών δεδομένων όπως σήματα ΗΚΓ, γονιδιωματικές αλληλουχίες και φασματοσκοπία [5]. Η λειτουργία της συνέλιξης περιλαμβάνει την ολίσθηση ενός φίλτρου (kernel) μήκους $K$ πάνω στο διάνυσμα εισόδου $x$. 

Για κάθε θέση $i$, υπολογίζεται το εσωτερικό γινόμενο
$$y[i] = (x * w)[i] = \sum_{k=0}^{K-1} x[i+k] \cdot w[k] + b$$
Στις προτεινόμενες αρχιτεκτονικές της παρούσας μελέτης, ένα 1D-CNN τοποθετείται \emph{πριν} το RNN, ως επίπεδο εξαγωγής χαρακτηριστικών (feature extraction): μαθαίνει να αναγνωρίζει τοπικά μοτίβα (συσχετίσεις μεταξύ γειτονικών τιμών στο διάνυσμα πρωτεϊνών, π.χ. ομάδες πρωτεϊνών που εκφράζονται μαζί), λειτουργεί ως φίλτρο εξομάλυνσης μειώνοντας τον υψίσυχνο θόρυβο της φασματομετρίας μάζας, και προσφέρει translation invariance (αναγνωρίζει ένα μοτίβο ανεξάρτητα από την ακριβή θέση του στην είσοδο).

\subsection{2.3 Μέθοδοι Συνόλου και Δέντρα Αποφάσεων}

Όπως αναφέρθηκε στην αρχή του κεφαλαίου, σε datasets με μικρό αριθμό δειγμάτων ($N < 100$) τα βαθιά νευρωνικά δίκτυα αντιμετωπίζουν σοβαρούς περιορισμούς. Η βιβλιογραφία υποδεικνύει ότι σε τέτοιες περιπτώσεις οι \textbf{Μέθοδοι Συνόλου} (Ensemble Methods) που βασίζονται σε δέντρα αποφάσεων αποτελούν μια αξιόπιστη εναλλακτική, καθώς παρέχουν ανθεκτικότητα στην υπερ-προσαρμογή, ερμηνευσιμότητα και υψηλή απόδοση χωρίς να απαιτούν μεγάλο πλήθος δειγμάτων [79], [80]. 

Τα \textbf{δέντρα αποφάσεων} (Decision Trees) είναι μη-παραμετρικοί αλγόριθμοι επιβλεπόμενης μάθησης που διαμερίζουν αναδρομικά τον χώρο των χαρακτηριστικών σε υποπεριοχές ομοιογενών κλάσεων. Σε κάθε εσωτερικό κόμβο, ο αλγόριθμος επιλέγει το χαρακτηριστικό $X_j$ και το κατώφλι $\theta$ που μεγιστοποιούν τον διαχωρισμό των κλάσεων. 

Η διαδικασία επαναλαμβάνεται αναδρομικά μέχρι να ικανοποιηθεί κάποιο κριτήριο τερματισμού (π.χ. ελάχιστος αριθμός δειγμάτων ανά φύλλο ή μέγιστο βάθος). Για την επιλογή του βέλτιστου διαχωρισμού χρησιμοποιούνται δύο κύρια κριτήρια: η \emph{Gini Impurity} $\text{Gini}(t) = 1 - \sum_{k=1}^{C} p_k^2$ (όπου $p_k$ είναι η αναλογία δειγμάτων της κλάσης $k$ στον κόμβο $t$), η οποία μετράει την πιθανότητα ένα τυχαία επιλεγμένο δείγμα να ταξινομηθεί λανθασμένα αν ταξινομηθεί σύμφωνα με την κατανομή των κλάσεων στον κόμβο, και η \emph{Information Gain (Entropy)} που βασίζεται στη θεωρία πληροφοριών του Shannon, με την εντροπία να ορίζεται ως $H(t) = -\sum_{k=1}^{C} p_k \log_2(p_k)$ και το κέρδος πληροφορίας από τη διαίρεση βάσει του χαρακτηριστικού $A$ ως $IG(t, A) = H(t) - \sum_{v} (|t_v|/|t|) H(t_v)$. 

Παρά την ερμηνευσιμότητά τους, τα μεμονωμένα δέντρα πάσχουν από \emph{υψηλή διακύμανση} (high variance): μικρές αλλαγές στα δεδομένα εκπαίδευσης μπορούν να οδηγήσουν σε εντελώς διαφορετικά δέντρα. Αυτός ο περιορισμός είναι ο λόγος ύπαρξης των μεθόδων συνόλου.

Ο Breiman (1996) εισήγαγε την τεχνική \textbf{Bagging} (Bootstrap Aggregating) ως μέθοδο μείωσης της διακύμανσης ασταθών εκτιμητών [79]. Η μέθοδος βασίζεται σε δύο θεμελιώδεις αρχές. Πρώτον, στο \emph{Bootstrap Sampling}: δημιουργία $B$ συνόλων εκπαίδευσης μέσω δειγματοληψίας με επανάθεση (sampling with replacement) από το αρχικό σύνολο μεγέθους $N$. Κάθε bootstrap sample περιέχει κατά μέσο όρο το $63.2\%$ των μοναδικών παρατηρήσεων του αρχικού συνόλου, όπως προκύπτει από τη σχέση $\lim_{N \to \infty} 1 - (1 - 1/N)^N = 1 - 1/e \approx 0.632$. Τα εναπομείναντα δείγματα, γνωστά ως \emph{Out-of-Bag (OOB) samples}, χρησιμοποιούνται ως εσωτερικό σύνολο επικύρωσης. Δεύτερον, στο \emph{Aggregation}: εκπαίδευση ενός μοντέλου βάσης (base learner) σε κάθε bootstrap sample και συνδυασμός των προβλέψεων μέσω πλειοψηφικής ψήφου για ταξινόμηση ή μέσου όρου για παλινδρόμηση,
$$\hat{y}_{\text{bag}}(x) = \text{mode}\{f_1(x), f_2(x), \ldots, f_B(x)\}$$
Για $B$ ανεξάρτητους εκτιμητές με διακύμανση $\sigma^2$, η διακύμανση του μέσου τους είναι $\sigma^2 / B$. Στην πράξη, οι εκτιμητές εμφανίζουν συσχέτιση $\rho$, οπότε η διακύμανση του ensemble δίνεται από
$$\text{Var}(\hat{y}_{\text{bag}}) = \rho \sigma^2 + \frac{1 - \rho}{B} \sigma^2$$
Ο δεύτερος όρος μειώνεται καθώς αυξάνεται το $B$, αλλά ο πρώτος παραμένει σταθερός. Η μείωση της συσχέτισης $\rho$ αποτελεί ακριβώς το κίνητρο για την επόμενη μέθοδο.

Το \textbf{Random Forest} (Breiman, 2001) επεκτείνει το Bagging με την εισαγωγή \emph{τυχαιότητας στην επιλογή χαρακτηριστικών} σε κάθε κόμβο διαχωρισμού [80]. Αντί να εξετάζονται όλα τα $P$ χαρακτηριστικά, σε κάθε split επιλέγεται ένα τυχαίο υποσύνολο μεγέθους $m$, όπου τυπικά $m = \lfloor \sqrt{P} \rfloor$ για classification και $m = \lfloor P/3 \rfloor$ για regression. Αυτός ο μηχανισμός μειώνει τη συσχέτιση μεταξύ των δέντρων: ακόμα και αν υπάρχει ένα πολύ ισχυρό χαρακτηριστικό (π.χ.~η MMP12 στο dataset μας), δεν θα χρησιμοποιείται σε κάθε δέντρο. Αυτό οδηγεί σε χαμηλότερη συνολική διακύμανση και καλύτερη γενίκευση. Ένα σημαντικό πλεονέκτημα του Random Forest, ιδιαίτερα στο πλαίσιο ανακάλυψης βιοδεικτών, είναι η δυνατότητα εκτίμησης της σημαντικότητας κάθε χαρακτηριστικού μέσω του δείκτη Mean Decrease in Impurity (MDI):
$$\text{Imp}(X_j) = \frac{1}{B} \sum_{b=1}^{B} \sum_{t \in T_b} \mathbb{I}(v(t) = j) \cdot p(t) \cdot \Delta\text{Gini}(t)$$
όπου $v(t)$ είναι το χαρακτηριστικό που χρησιμοποιείται στον κόμβο $t$, $p(t)$ η αναλογία δειγμάτων που φτάνουν στον κόμβο $t$, και $\Delta\text{Gini}(t)$ η μείωση ακαθαρσίας Gini λόγω του διαχωρισμού. Η μετρική αυτή παρέχει μια εναλλακτική, ανεξάρτητη μέθοδο ταυτοποίησης βιοδεικτών η οποία μπορεί να συγκριθεί με την ANOVA $F$-scores ανάλυση. Παράλληλα, το \emph{OOB error} προσφέρει εκτίμηση του σφάλματος γενίκευσης χωρίς ξεχωριστό validation set, ιδιαίτερα χρήσιμο σε datasets μικρού μεγέθους [83].

Πέρα από το κλασικό Random Forest, αναπτύχθηκαν και πιο σύνθετες μέθοδοι συνόλου που αυξάνουν περαιτέρω την ποικιλομορφία μεταξύ των δέντρων. Το \textbf{Rotation Forest} (Rodriguez et al., 2006) διαμερίζει τυχαία τα χαρακτηριστικά σε ομάδες, εφαρμόζει PCA σε κάθε ομάδα, και συνενώνει τα principal components ως νέα χαρακτηριστικά. Κάθε δέντρο, έτσι, ``βλέπει'' τα δεδομένα από μια εντελώς διαφορετική οπτική γωνία. Το \textbf{Random Subspace} εκπαιδεύει κάθε δέντρο σε ένα τυχαίο υποσύνολο όλων των χαρακτηριστικών (όχι μόνο σε κάθε κόμβο όπως το RF), αυξάνοντας τη συνολική ποικιλομορφία.

Μια ιδιαίτερα ενδιαφέρουσα πρόσφατη εξέλιξη είναι το \textbf{Deep Forest} (Zhou \& Feng, 2017), επίσης γνωστό ως \emph{gcForest}. Οι συγγραφείς παρατήρησαν ότι η επιτυχία της Βαθιάς Μάθησης βασίζεται σε τρεις βασικές αρχές: (α) επεξεργασία κατά στρώματα (layer-by-layer processing), (β) μετασχηματισμός χαρακτηριστικών εντός κάθε στρώματος (in-model feature transformation), και (γ) επαρκής πολυπλοκότητα μοντέλου. Αυτές οι αρχές δεν είναι, όμως, αποκλειστικά συνδεδεμένες με τα νευρωνικά δίκτυα. Το Deep Forest ενσωματώνει τις ίδιες αρχές χρησιμοποιώντας αποκλειστικά δέντρα αποφάσεων και μεθόδους συνόλου, χωρίς backpropagation [81], [82]. Τα κύρια πλεονεκτήματα είναι σημαντικά λιγότερες υπερπαράμετροι (δεν απαιτείται ρύθμιση learning rate, batch size ή αρχιτεκτονικής), αυτόματη ρύθμιση βάθους (μέσω early stopping βάσει cross-validation), και εξαιρετική απόδοση σε μικρά datasets (σχεδιασμένο για $N < 1000$). Η αρχιτεκτονική αποτελείται από μια σειρά επιπέδων (cascade) όπου κάθε επίπεδο περιέχει πολλαπλούς και διαφορετικούς τύπους forests, τυπικά Random Forests, Completely Random Forests (στα οποία τόσο τα χαρακτηριστικά όσο και τα κατώφλια επιλέγονται τυχαία) και Extra Trees [84]. Έστω $x \in \mathbb{R}^d$ το αρχικό διάνυσμα εισόδου και $C$ ο αριθμός κλάσεων. Σε κάθε επίπεδο $l$ του cascade, κάθε forest $f$ ($f = 1, \ldots, F$) παράγει ένα class probability vector $p_f^{(l)} \in \mathbb{R}^C$ μέσω εσωτερικού $k$-fold cross-validation στο σύνολο εκπαίδευσης (αυτό αποτρέπει information leakage μεταξύ επιπέδων). Τα probability vectors συνενώνονται με το αρχικό διάνυσμα εισόδου, σχηματίζοντας ένα augmented feature vector
$$x^{(l+1)} = \left[ x \; ; \; \underset{f=1}{\overset{F}{\text{concat}}} \; p_f^{(l)} \right] \in \mathbb{R}^{d + F \cdot C}$$
και η διαδικασία επαναλαμβάνεται. Η προσθήκη νέων επιπέδων σταματά αυτόματα όταν η απόδοση στο validation set δεν βελτιώνεται, αποφεύγοντας την υπερ-προσαρμογή. Στο τελικό επίπεδο $L$, η πρόβλεψη υπολογίζεται ως ο μέσος όρος των class probability vectors:
$$\hat{y} = \arg\max_c \frac{1}{F} \sum_{f=1}^{F} p_{f,c}^{(L)}$$

Η επιλογή μεθοδολογίας, νευρωνικά δίκτυα έναντι μεθόδων συνόλου, εξαρτάται κρίσιμα από τα χαρακτηριστικά του dataset. Συνοπτικά, τα tree-based μοντέλα απαιτούν τυπικά 50--100 δείγματα, παρουσιάζουν χαμηλό κίνδυνο overfitting, δεν χρειάζονται κανονικοποίηση δεδομένων, προσφέρουν υψηλή ερμηνευσιμότητα μέσω Feature Importance, αλλά μοντελοποιούν χρονικές δομές μόνο χειροκίνητα (μέσω feature engineering). Τα νευρωνικά δίκτυα (RNN/LSTM), αντιθέτως, τυπικά απαιτούν $\sim 1000+$ δείγματα, έχουν υψηλό κίνδυνο overfitting, απαιτούν κανονικοποίηση, λειτουργούν ως black box (χαμηλή ερμηνευσιμότητα χωρίς attention), αλλά μοντελοποιούν χρονικές δομές εγγενώς. Πρόσφατες μελέτες [85] έχουν αναδείξει ότι σε ένα ευρύ φάσμα δομημένων (tabular) datasets με $N < 10{,}000$, τα tree-based models (ιδίως XGBoost και Random Forest) υπερτερούν σταθερά των νευρωνικών δικτύων, εύρημα ιδιαίτερα σχετικό με τη φύση του πρωτεομικού dataset της παρούσας μελέτης.
